% Section 6, part 8.
\subsection{Theorem V: logarithmic legs}
\label{sec:mod-V}

Here the model curves form a family in $v\neq 0$; $\partial_v|_x$ is the derivative in $v$ at fixed $x$, computed away from the branch points through the inverse branches $p(x,v)$, and $\partial_v|_p$ the derivative at fixed $p$.

\begin{lem}[Fixed $p$ versus fixed $x$]
\label{lem:mod-fixedx}
$\partial_vp|_x=-p/(p^5-v)$. For every family of meromorphic forms, $\partial_v|_x=\partial_v|_p-\mathrm{Lie}_{\cX}$ in each variable, where $\cX:=\frac{p}{p^5-v}\frac{\dd}{\dd p}$. Moreover $\partial_v(y\,\dd x)|_x=-\dd p/p=\int_{q\in\mathfrak r}\omega_{0,2}(\cdot,q)$.
\end{lem}

\begin{proof}
$\partial_vp|_x=-(\partial_vx|_p)/(\partial_px)=-(1/p)\big/\big((p^5-v)/p^2\big)$. For $\omega=f(p,v)\,\dd p$ the chain rule gives $\partial_v|_x\omega=(\partial_vf+f_p\,\partial_vp|_x)\dd p+f\,\dd(\partial_vp|_x)=\partial_v|_p\omega-\mathrm{Lie}_{\cX}\omega$, in each variable. Since $y=p$, $\partial_v(y\,\dd x)|_x=(\partial_vp|_x)\,\dd x=-\dd p/p$, and $\int_{q\in\mathfrak r}\frac{\dd p\,\dd q}{(p-q)^2}=\dd p\big[\frac{1}{p-q}\big]_{q=0}^{q=\infty}=-\frac{\dd p}{p}$.
\end{proof}

\begin{lem}[Integration by parts]
\label{lem:mod-ibp}
Let $w$ be a meromorphic $1$-form on $\bbP^1$ with poles only at the branch points, zero residues, holomorphic at $p=0$ and $O(\dd p/p^2)$ at $p=\infty$ (for instance $\omega_{h,n}$ in one slot, $2h-2+n>0$).
\begin{enumerate}
\item For $\phi\in\bbQ[v][p]$ let $(\cX\dd\phi)_+\in\bbQ[v][p]$ be the quotient in the division $p\,\phi'(p)=(\cX\dd\phi)_+\,(p^5-v)+r$, $\deg r<5$; for $\phi=p^m$, $(\cX\dd\phi)_+=m\sum_{j\geq 0,\ m-5-5j\geq 0}v^jp^{m-5-5j}$. Then $(2\pi\ri)^{-1}\oint_{|p|=R}\phi\,\mathrm{Lie}_{\cX}w$ equals $-\langle(\cX\dd\phi)_+,w\rangle$.
\item $\int_{\mathfrak r}\mathrm{Lie}_{\cX}w=0$.
\end{enumerate}
\end{lem}

\begin{proof}
(i) $\mathrm{Lie}_{\cX}(\phi w)=\dd(\iota_{\cX}(\phi w))$ has zero circle integral, and $\mathrm{Lie}_{\cX}(\phi w)=(\cX\phi)\,w+\phi\,\mathrm{Lie}_{\cX}w$ with $\cX\phi=(\cX\dd\phi)_++r/(p^5-v)$. The form $\frac{r}{p^5-v}w$ has all its poles inside the circle and is $O(\dd p/p^3)$ at $\infty$, so its circle integral vanishes; and the circle pairing of a polynomial is its residue pairing (\cref{def:mod-pair}). For $\phi=p^m$, expand $mp^m/(p^5-v)$ for $|p|^5>|v|$. (ii) $\int_{\mathfrak r}\mathrm{Lie}_{\cX}w=[\iota_{\cX}w]_0^\infty=0$: $\cX$ vanishes at $0$, where $w$ is holomorphic, and $\cX=O(p^{-4})\,\dd/\dd p$, $w=O(\dd p/p^2)$ at $\infty$.
\end{proof}

\begin{thm}[The variational formula, in the form used]
\label{thm:mod-var}
Let $J=(p_1,\dots,p_n)$ be points off the branch points.
\begin{enumerate}
\item For $n\geq 1$ and $h+n>1$: $\partial_v|_x\,\omega_{h,n}(J)=\int_{q\in\mathfrak r}\omega_{h,n+1}(J,q)$.
\item For $h\geq 2$: $\partial_v\FTR_h=\int_{\mathfrak r}\omega_{h,1}$.
\end{enumerate}
\end{thm}

\begin{proof}
This is \cite[Thm.~5.1]{Eynard:2007kz}, (5-14) for $g+k>1$ and (5-15) for $g\geq 1$, for the temperature variation of \cite[\S5.3, (5-27)--(5-30)]{Eynard:2007kz}, translated by \cref{conv:mod-kernel}. There a variation at fixed $x$ is $\delta_\Omega y|_x\,\dd x=-\Omega$ \cite[(5-1)]{Eynard:2007kz} with $\Omega(p)=\int_{\partial\Omega}B(p,q)\Lambda(q)$ \cite[(5-13)]{Eynard:2007kz}, and (5-14), (5-15) read $\delta_\Omega W^{(h)}_n(J)|_{x}=\int_{\partial\Omega}W^{(h)}_{n+1}(J,q)\Lambda(q)$, $\delta_\Omega F^{(h)}=-\int_{\partial\Omega}W^{(h)}_1\Lambda$. By \cref{lem:mod-fixedx}, $\Omega=-\partial_v(y\,\dd x)|_x=-\int_{\mathfrak r}\omega_{0,2}(\cdot,q)$: the path $\partial\Omega=\mathfrak r$ joins the poles $0$, $\infty$ of $y\,\dd x$, with $\Lambda=-1$. With $\omega_{h,n}=(-1)^nW^{(h)}_n$ and $\FTR_h=-F^{(h)}$ this gives $\partial_v|_x\omega_{h,n}=(-1)^{n}(-1)^{n+1}(-1)\int_{\mathfrak r}\omega_{h,n+1}=\int_{\mathfrak r}\omega_{h,n+1}$ and $\partial_v\FTR_h=\int_{\mathfrak r}W^{(h)}_1\Lambda=\int_{\mathfrak r}\omega_{h,1}$. The hypotheses hold: the branch points are simple with $\dd y\neq 0$; the variation is at fixed $x$ with the Bergman kernel of the fixed $\bbP^1$; $\Omega/\dd y=1/p$ is regular at the branch points \cite[(5-4)]{Eynard:2007kz}; the path avoids them, without hard edges \cite[(5-13), footnote~8]{Eynard:2007kz}; the integrals converge ($\omega_{h,n+1}$ is holomorphic at $0$, $\infty$); and residues at the moving branch points commute with $\partial_v|_x$ (fixed circles in the $x$-plane).
\end{proof}

\begin{lem}[Genus zero]
\label{lem:mod-genus0}
Put $e_{b,k}(p):=\dd p/(p-\varsigma_b)^{k+1}$ and $r_b:=1/(5\varsigma_b^3)$. Then:
\begin{enumerate}
\item $\partial_v|_x\,\omega_{0,2}(p_1,p_2)=\sum_br_b\,e_{b,1}(p_1)\,e_{b,1}(p_2)$;
\item $\omega_{0,3}(p_0,p_1,p_2)=-\sum_be_{b,1}(p_0)e_{b,1}(p_1)e_{b,1}(p_2)/x''(\varsigma_b)$, the negative of \cite[Thm.~4.1 (4-20)]{Eynard:2007kz};
\item $\int_{q\in\mathfrak r}\omega_{0,3}(p_1,p_2,q)=\partial_v|_x\,\omega_{0,2}(p_1,p_2)$;
\item for $\phi_1,\phi_2,\phi_3\in\bbQ[v][p]$,
\[
\begin{aligned}
\langle\phi_1\phi_2\phi_3\rangle^{\mathrm{TR}}_0&=-\tfrac15\textstyle\sum_b\varsigma_b^{-2}\prod_i\phi_i'(\varsigma_b), &\langle\phi_1\phi_2,-\log p\rangle^{\mathrm{TR}}_0&=\tfrac15\textstyle\sum_b\varsigma_b^{-3}\phi_1'(\varsigma_b)\phi_2'(\varsigma_b),\\
\langle\phi_1,(-\log p)^2\rangle^{\mathrm{TR}}_0&=-\tfrac15\textstyle\sum_b\varsigma_b^{-4}\phi_1'(\varsigma_b), &\langle(-\log p)^3\rangle^{\mathrm{TR}}_0&=1/v,
\end{aligned}
\]
and every pairing of $\omega_{0,3}$ with at least one leg in $\bbQ[v][p]$ lies in $\bbQ[v]$.
\end{enumerate}
\end{lem}

\begin{proof}
(i) $\omega_{0,2}$ does not depend on $v$ at fixed $p$, so by \cref{lem:mod-fixedx} $\partial_v|_x\omega_{0,2}=-(\mathrm{Lie}_{\cX,1}+\mathrm{Lie}_{\cX,2})\,\omega_{0,2}$. For a function $g$ and the field $g\,\dd/\dd p$ in both variables,
\[
(\mathrm{Lie}_{1}+\mathrm{Lie}_{2})\frac{\dd p_1\dd p_2}{(p_1-p_2)^2}=\Big(\frac{g'(p_1)+g'(p_2)}{(p_1-p_2)^2}-\frac{2(g(p_1)-g(p_2))}{(p_1-p_2)^3}\Big)\dd p_1\dd p_2=\dd_1\dd_2\frac{g(p_1)-g(p_2)}{p_1-p_2},
\]
and with $g=p/(p^5-v)=\sum_br_b/(p-\varsigma_b)$, $(g(p_1)-g(p_2))/(p_1-p_2)=-\sum_br_b/((p_1-\varsigma_b)(p_2-\varsigma_b))$.
(ii) Only the two terms of ${\sum}'$ with $I_1=\{p_1\}$, $\{p_2\}$ occur. With $t:=q-\varsigma_b$, $\iota_bq=\varsigma_b-t+O(t^2)$, $\dd x(q)=x''(\varsigma_b)t(1+O(t))\dd t$, $\int_{\iota_bq}^q\omega_{0,2}(p_0,\cdot)=2t\,e_{b,1}(p_0)(1+O(t))$ and $y(q)-y(\iota_bq)=2t(1+O(t))$, so $K_b(p_0,q)=\frac{e_{b,1}(p_0)}{2x''(\varsigma_b)t\,\dd t}(1+O(t))$ and the bracket is $-2\,e_{b,1}(p_1)e_{b,1}(p_2)\,\dd t^2(1+O(t))$; the residue is the claim.
(iii) $x''(\varsigma_b)=5\varsigma_b^2$ and $\int_{\mathfrak r}e_{b,1}=\big[-1/(p-\varsigma_b)\big]_0^\infty=-1/\varsigma_b$, so (ii) gives $\sum_br_b\,e_{b,1}(p_1)e_{b,1}(p_2)$, which is (i).
(iv) $\Res_{\varsigma_b}\phi\,e_{b,1}=\phi'(\varsigma_b)$ and $\langle-\log p,e_{b,1}\rangle=-1/\varsigma_b$; inserting into (ii) gives the four formulas ($\frac15\sum_b\varsigma_b^{-5}=1/v$). Constants pair to zero, so with an L-leg we may take $\phi_i\in p\,\bbQ[v][p]$; then only powers $\varsigma_b^N$, $N\geq -4$, occur, and $\sum_b\varsigma_b^N=5v^{N/5}$ if $5\mid N$, $0$ otherwise.
\end{proof}

\begin{prop}[The recursion (V1)]
\label{prop:mod-V1}
Let $h\geq 0$, $a,k\geq 0$, $n=a+k\geq 1$, $2h-2+n>0$, and $\phi_1,\dots,\phi_a\in\bbQ[v][p]$. Then
\begin{equation}
\label{eq:mod-V1}
\begin{aligned}
\langle\phi_1\cdots\phi_a,(-\log p)^{k+1}\rangle^{\mathrm{TR}}_h={}&\partial_v\langle\phi_1\cdots\phi_a,(-\log p)^k\rangle^{\mathrm{TR}}_h-\sum_{i=1}^a\langle\cdots,\partial_v\phi_i,\cdots,(-\log p)^k\rangle^{\mathrm{TR}}_h\\
&+\sum_{i=1}^a\langle\cdots,(\cX\dd\phi_i)_+,\cdots,(-\log p)^k\rangle^{\mathrm{TR}}_h ,
\end{aligned}
\end{equation}
where $\partial_v\phi_i$ is the derivative of the coefficients and the dots stand for the other legs.
\end{prop}

\begin{proof}
Fix $v_0\neq 0$; for $v$ near $v_0$ take $R>\max_b|\varsigma_b(v)|$ and a ray $\mathfrak r$ avoiding all $\varsigma_b(v)$. Then $C(v):=\langle\phi_1\cdots\phi_a,(-\log p)^k\rangle^{\mathrm{TR}}_h=(2\pi\ri)^{-a}\oint_{|p_i|=R}\int_{\mathfrak r^k}\prod_i\phi_i(p_i)\,\omega_{h,n}$ with contours independent of $v$, and the integrand is holomorphic in $v$, bounded on the rays by a constant times $|\dd q|/(1+|q|)^2$ uniformly near $v_0$. Differentiating under the integrals, $\partial_vC=\sum_i\langle\cdots,\partial_v\phi_i,\cdots\rangle^{\mathrm{TR}}_h$ plus the integral of $\prod_i\phi_i\,\partial_v|_p\omega_{h,n}$. By \cref{lem:mod-fixedx} and \cref{thm:mod-var}~(i), $\partial_v|_p\omega_{h,n}=\int_{q\in\mathfrak r}\omega_{h,n+1}(\cdot,q)+\sum_{l=1}^n\mathrm{Lie}_{\cX,l}\,\omega_{h,n}$. The first term gives the pairing with one more log leg (Fubini); a Lie term gives $-\langle\cdots,(\cX\dd\phi_i)_+,\cdots\rangle^{\mathrm{TR}}_h$ in an L-slot and $0$ in a log slot (\cref{lem:mod-ibp}). Rearranging gives \eqref{eq:mod-V1}.
\end{proof}

% Section 6, part 9.
Let $c$ be a branch point, $t:=p-c$, and $\chi$ (a local letter) a local coordinate at $c$ with $x=x(c)+\lambda\chi^2$, $\lambda\neq 0$, so that $\iota_c\colon\chi\mapsto-\chi$. Write $y=\sum_ky_k\chi^k$ and, near $(c,c)$, $\omega_{0,2}=\big[(\chi_1-\chi_2)^{-2}+\mathrm{reg}(\chi_1,\chi_2)\big]\dd\chi_1\dd\chi_2$ with $\mathrm{reg}$ holomorphic and symmetric; then $\mathrm{reg}_{00}:=\mathrm{reg}(0,0)=\frac16\{t;\chi\}(0)$, where $\{t;\chi\}=t'''/t'-\frac32(t''/t')^2$ is the Schwarzian derivative of $t$ as a function of $\chi$. For a $1$-form $\eta=\sum_m\eta_m\chi^m\dd\chi$ holomorphic at $c$ put
\begin{equation}
\label{eq:mod-loc}
\mathrm{Loc}_c[\eta]:=-\frac{1}{4\lambda y_1}\Big[\mathrm{reg}_{00}\,\eta_0+\frac{\eta_2}{12}-\frac{y_3}{4y_1}\,\eta_0\Big].
\end{equation}

\begin{lem}[A local formula for $\omega_{1,1}$]
\label{lem:mod-w11}
\begin{enumerate}
\item $\omega_{1,1}(p_0)=\sum_c\mathrm{Loc}_c[\omega_{0,2}(p_0,\cdot)]$ for $p_0$ off the branch points, $\omega_{0,2}(p_0,\cdot)$ being a form in the second slot.
\item $\omega_{1,1}=\dd G$ with $G$ rational, with poles only at the branch points and $G(\infty)=0$; so $\omega_{1,1}$ has zero residues, is holomorphic at $p=0$ and is $O(\dd p/p^2)$ at $\infty$.
\item $\langle-\log p\rangle^{\mathrm{TR}}_1=\int_{\mathfrak r}\omega_{1,1}=\sum_c\mathrm{Loc}_c[-\dd q/q]$.
\end{enumerate}
\end{lem}

\begin{proof}
(i) By \eqref{eq:mod-kernel},
\[
\omega_{1,1}(p_0)=\sum_c\Res_{q\to c}K_c(p_0,q)\,\omega_{0,2}(q,\iota_cq).
\]
Expand $\omega_{0,2}(p_0,q')=\dd p_0\sum_mB_m(p_0)\chi^m\dd\chi$ for $q'$ near $c$ with coordinate $\chi$. Then $\int_{\iota_cq}^q\omega_{0,2}(p_0,\cdot)=2\,\dd p_0\sum_{m\text{ even}}B_m\chi^{m+1}/(m+1)$, $y(q)-y(\iota_cq)=2\sum_{k\text{ odd}}y_k\chi^k$ and $\dd x=2\lambda\chi\,\dd\chi$, so
\[
K_c(p_0,q)=\frac{\dd p_0}{\dd\chi}\,\frac{1}{4\lambda y_1}\Big[B_0\chi^{-1}+\Big(\frac{B_2}{3}-B_0\frac{y_3}{y_1}\Big)\chi+O(\chi^3)\Big],
\]
and $\omega_{0,2}(q,\iota_cq)=-\big[(4\chi^2)^{-1}+\mathrm{reg}_{00}+O(\chi^2)\big]\dd\chi^2$ ($\dd(-\chi)=-\dd\chi$, and $\mathrm{reg}(\chi,-\chi)=\mathrm{reg}_{00}+O(\chi^2)$ by symmetry). The coefficient of $\chi^{-1}\dd\chi$ in the product is $-\frac{1}{4\lambda y_1}\big[B_0\,\mathrm{reg}_{00}+\frac14\big(\frac{B_2}{3}-B_0\frac{y_3}{y_1}\big)\big]\dd p_0=\mathrm{Loc}_c[\omega_{0,2}(p_0,\cdot)]$.
(ii) $\omega_{0,2}(p_0,q)=\dd_{p_0}\big[\dd q/(q-p_0)\big]$, and near $c$, $\dd q/(q-p_0)=\sum_mg_m(p_0)\chi^m\dd\chi$ with $g_m$ rational, with poles only at $c$ and $g_m(\infty)=0$. So $B_m\dd p_0=\dd g_m$, and by (i) $\omega_{1,1}=\dd G$ with $G$ a combination of the $g_0,g_2$ at the five branch points. (iii) The residues in (i) are integrals over small circles avoiding $\mathfrak r$; the ray integral converges by (ii), and as $\omega_{0,2}(p_0,q)=O(p_0^{-2})\dd p_0$ uniformly for $q$ on the circles, it may be taken inside: $\int_{p_0\in\mathfrak r}\omega_{0,2}(p_0,q)=-\dd q/q$, and $\mathrm{Loc}_c$ is linear in finitely many Taylor coefficients. Finally, by (i) and \eqref{eq:mod-kernel}, $\mathrm{Loc}_c[\eta]=\Res_{q\to c}\frac12\frac{\int_{\iota_cq}^q\eta}{(y(q)-y(\iota_cq))\dd x(q)}\,\omega_{0,2}(q,\iota_cq)$ for $\eta$ holomorphic at $c$; this form does not depend on the coordinate $\chi$.
\end{proof}

\begin{thm}[Theorem V]
\label{thm:mod-V}
Let $h\geq 0$, $a,k\geq 0$, $n=a+k\geq 1$, $2h-2+n>0$, and $\phi_1,\dots,\phi_a\in\bbQ[v][p]$.
\begin{enumerate}
\item If $a\geq 1$, then $\langle\phi_1\cdots\phi_a,(-\log p)^k\rangle^{\mathrm{TR}}_h\in\bbQ[v]$.
\item If $a=0$, then $\langle(-\log p)^k\rangle^{\mathrm{TR}}_h=(\FTR_h)^{(k)}(v)$, with $\FTR_h$ of \cref{thm:mod-gap} for $h\geq 2$, $\FTR_1:=-\frac1{12}\log v$ (here $k\geq 1$) and $\FTR_0:=\frac12v^2\log v$ (here $k\geq 3$).
\item The recursion \eqref{eq:mod-V1} holds.
\end{enumerate}
\end{thm}

\begin{proof}
(iii) is \cref{prop:mod-V1}. (i) Fix $h$ and $a\geq 1$ and induct on $k$. The first admissible case is $k=0$ if $2h-2+a>0$ (\cref{thm:mod-Lreg}), and otherwise $\langle\phi_1,(-\log p)^2\rangle^{\mathrm{TR}}_0$ or $\langle\phi_1\phi_2,-\log p\rangle^{\mathrm{TR}}_0$ (\cref{lem:mod-genus0}~(iv)). In the step, the right side of \eqref{eq:mod-V1} consists of a $v$-derivative of, and pairings with, $a$ L-legs in $\bbQ[v][p]$ ($p^5-v$ is monic) and $k$ log legs; all lie in $\bbQ[v]$ by induction.

(ii) For $a=0$, \eqref{eq:mod-V1} reads $\langle(-\log p)^{k+1}\rangle^{\mathrm{TR}}_h=\partial_v\langle(-\log p)^k\rangle^{\mathrm{TR}}_h$, so it suffices to treat the first admissible $k$. \emph{$h\geq 2$, $k=1$.} As in the proof of \cref{prop:mod-period}~(iii), $(2-2h)\FTR_h=v\langle-\log p\rangle^{\mathrm{TR}}_h$, and $\FTR_h$ is a multiple of $v^{2-2h}$ (\cref{thm:mod-gap}), so $\langle-\log p\rangle^{\mathrm{TR}}_h=(\FTR_h)'(v)$. \emph{$h=0$, $k=3$.} \Cref{lem:mod-genus0}~(iv). \emph{$h=1$, $k=1$.} We show $\langle-\log p\rangle^{\mathrm{TR}}_1=-1/(12v)$.

\emph{Reduction to $v=1$.} By \cref{lem:mod-scale}~(ii) with $k=\varsigma$, $m_\varsigma^*\omega^{(v)}_{1,1}=v^{-1}\omega^{(1)}_{1,1}$, and rays go to rays, so $\langle-\log p\rangle^{\mathrm{TR}}_1(v)=v^{-1}\langle-\log p\rangle^{\mathrm{TR}}_1(1)$. Let $v=1$, $\varsigma=1$, $\varsigma_b=\zeta^{-b}$.

\emph{Five equal contributions.} $\rho_b(p):=\zeta^{-b}p$ sends $1$ to $\varsigma_b$, with $x\circ\rho_b=\zeta^bx$, $y\circ\rho_b=\zeta^{-b}y$, $\rho_b^*\omega_{0,2}=\omega_{0,2}$, $\rho_b^*(\dd q/q)=\dd q/q$; hence $\iota_b\circ\rho_b=\rho_b\circ\iota_0$ and $(y(q)-y(\iota q))\,\dd x(q)$ is unchanged. By the coordinate-free form of $\mathrm{Loc}$, $\mathrm{Loc}_{\varsigma_b}[-\dd q/q]=\mathrm{Loc}_1[-\dd q/q]$.

\emph{The contribution at $c=1$.} With $t=p-1$, $x(1+t)=\frac54+\frac52t^2+\frac54t^4-t^5+O(t^6)$, so $\lambda=5/2$, $\chi=t+\frac{t^3}4-\frac{t^4}5+O(t^5)$ and $t=\chi-\frac{\chi^3}4+\frac{\chi^4}5+O(\chi^5)$. Hence $y_1=1$, $y_2=0$, $y_3=-\frac14$; $\{t;\chi\}(0)=-\frac32$ ($t'''(0)=-\frac32$, $t''(0)=0$), so $\mathrm{reg}_{00}=-\frac14$; and $-\dd q/q=-\big(1-\chi+\frac14\chi^2+O(\chi^3)\big)\dd\chi$, so $\eta_0=-1$, $\eta_2=-\frac14$. By \eqref{eq:mod-loc}, $\mathrm{Loc}_1[-\dd q/q]=-\frac1{10}\big[\frac14-\frac1{48}-\frac1{16}\big]=-\frac1{60}$. By \cref{lem:mod-w11}~(iii) the total at $v=1$ is $5\cdot(-\frac1{60})=-\frac1{12}$. So $\langle-\log p\rangle^{\mathrm{TR}}_1=-1/(12v)=(\FTR_1)'(v)$.
\end{proof}

% Section 6, part 10.
\subsection{The Model Lemma on \texorpdfstring{$\sK$}{K}}
\label{sec:mod-MLK}

Recall $\ML{h}$ from \cref{def:red-ML}, with $\Fm_h$ the free energy of \cref{cor:mod-trF} for $h\geq 2$, $\Fm_1=-\frac1{12}\log v$ and $\Fm_0=-\frac52v^2\log v$ (only their derivatives of order $\geq 1$, respectively $\geq 3$, enter). Then $\Fm_h=(-5)^{1-h}\FTR_h$ for every $h\geq 0$, with $\FTR_0,\FTR_1$ as in \cref{thm:mod-V}~(ii).

\begin{thm}[The Model Lemma on $\sK$]
\label{thm:mod-MLK}
$\ML{h}$ holds for every $h\geq 0$.
\end{thm}

\begin{proof}
Part (ii) of $\ML{h}$ is \cref{thm:mod-gap}. For part (i), let $X_i\in\sK\cap V_{w_i}$, $i=1,\dots,n$.

\emph{Step 1 (the legs of $\sK$).} By \cref{cor:mod-Ksplit}, $X_i=Y_i+\Pi_0(X_i)H$ with $Y_i\in\sKo\cap V_{w_i}$ and $\Pi_0(X_i)\in\bbQ v^{w_i/5}$, a monomial with a non-negative exponent. By \cref{prop:mod-LD} ($q_{-\log p}=H$, and the TR leg of $Y_i$ is $-\Lambda(Y_i)$), $X_i=q_{\ell_i}$ for $\ell_i:=-\Lambda(Y_i)+\Pi_0(X_i)\,(-\log p)$. By \cref{thm:mod-lattice}~(b) and \cref{lem:mod-classL}~(i), $-\Lambda(Y_i)=\phi_i+f_i(x)$ with $\phi_i\in\bbQ[v][p]$ and $f_i\in\bbQ[v][x]$.

\emph{Step 2 (expansion).} By \cref{thm:mod-tr}, multilinearity and \cref{lem:mod-xfun}, which removes the components $f_i(x)$ in every slot,
\[
\langle X_1\cdots X_n\rangle_h=(-5)^{1-h}\sum_{S\subseteq\{1,\dots,n\}}\prod_{i\notin S}\Pi_0(X_i)\,\big\langle(\phi_i)_{i\in S},(-\log p)^{n-|S|}\big\rangle^{\mathrm{TR}}_h .
\]
Every term with $S\neq\emptyset$ lies in $\bbQ[v]$, by \cref{thm:mod-Lreg} if $S=\{1,\dots,n\}$ and by \cref{thm:mod-V}~(i) otherwise.

\emph{Step 3 (the pure logarithmic term).} By \cref{thm:mod-V}~(ii) and $\Fm_h=(-5)^{1-h}\FTR_h$, the term $S=\emptyset$ is $(-5)^{1-h}\langle(-\log p)^n\rangle^{\mathrm{TR}}_h\prod_i\Pi_0(X_i)=(\Fm_h)^{(n)}\prod_i\Pi_0(X_i)$.
\end{proof}

\begin{cor}[The one-point law on $\sK$]
\label{cor:mod-onepoint}
Let $h\geq 1$ and $X\in\sK\cap V_w$. If $\langle X\rangle_h$ is singular, that is, if its $v$-exponent is negative, then $\langle X\rangle_h=\Pi_0(X)\,(\Fm_h)'(v)$.
\end{cor}

\begin{proof}
By \cref{thm:mod-MLK} with $n=1$, $\langle X\rangle_h-(\Fm_h)'\Pi_0(X)\in\bbQ[v]$. By \cref{lem:set-modhom} both terms are multiples of $v^e$, $e=1-2h+w/5<0$ (as $\Pi_0(X)\in\bbQ v^{w/5}$), so the difference is $0$.
\end{proof}

\begin{rmk}[Off $\sK$]
\label{rmk:mod-offK}
The restriction to $\sK$ cannot be dropped: for insertions with a component along some $\psi^kH$, $k\geq 1$, the singular part of a one-point correlator is in general not $\Pi_0(X)(\Fm_h)'$, nor are those of two-point correlators given by $\Fm_h(V(q))$ for a genus-independent $V$; both facts are known only by computation (\cref{sec:evi-model}) and are not used.
\end{rmk}


